% Abelian-square densities under bounded local avoidance
% v0.6.1 -- transcribed by hand from MANUSCRIPT_DRAFT_v0.5.md (authoritative text).
% No automated Markdown-to-TeX conversion was used.
\documentclass[12pt]{article}
\usepackage[utf8]{inputenc}
\usepackage[T1]{fontenc}
\usepackage{amsmath,amssymb,amsthm,booktabs,hyperref,graphicx}

\newtheorem{theorem}{Theorem}[section]
\newtheorem{lemma}[theorem]{Lemma}
\newtheorem{proposition}[theorem]{Proposition}
\newtheorem{corollary}[theorem]{Corollary}
\theoremstyle{definition}
\newtheorem{remark}[theorem]{Remark}

\title{Abelian-square densities under bounded local avoidance:\\ a non-monotone ternary family}
\author{[to be inserted]}
\date{2026-08-24}

\begin{document}
\maketitle

\begin{abstract}
Let \(L_h\) be the ternary words containing no Abelian square of half-length \(K\in\{2,\ldots,h\}\); half-length \(1\) is permitted, so \(aa,bb,cc\) remain allowed. We ask how often a word of \(L_h\) is \emph{itself} an Abelian square. Let \(R_n^{(h)}\) be the fraction of words in \(L_h\cap\Sigma^{2n}\) whose two halves have equal Parikh vector. For each \(h\in\{2,\ldots,6\}\) we prove
\[
R_n^{(h)}\sim C_h/n,
\qquad
C_h=\frac{1}{2\sqrt3\,\pi\,a_h},
\]
where \(a_h\) is the asymptotic variance rate of one letter count under the measure of maximal entropy on \(L_h\). The proof folds the event into a rank-two additive functional of a finite Markov chain, certifies a unimodular product lattice by explicit equal-length cycles, and shows that transient window states cancel exactly between numerator and denominator.

Every \(C_h\) exceeds the free ternary value \(3\sqrt3/(4\pi)\): within this family, bounded local avoidance raises the leading density of long Abelian squares. The dependence on \(h\) is not monotone---the effect weakens for \(h=2,3,4\) and grows at \(h=5,6\). The first increase occurs at the same cutoff as the first loss of recurrent states, and a \(K=5\) ablation identifies a single half-Parikh transition class associated with both.
\end{abstract}

\section{Introduction}
\label{sec:intro}

An Abelian square is a word \(UV\) with equal-length consecutive halves satisfying
\[
\Psi(U)=\Psi(V),
\]
where
\[
\Psi(w)=\bigl(|w|_a,|w|_b,|w|_c\bigr)
\]
is the Parikh vector.

Fix the ternary alphabet
\[
\Sigma=\{a,b,c\}.
\]
For \(h\ge2\), define
\[
L_h
=
\left\{
w\in\Sigma^*:
w\text{ contains no Abelian square of half-length }K\in\{2,\ldots,h\}
\right\}.
\]
The half-length \(K=1\) is deliberately excluded from the forbidden set: this matches the M\"akel\"a formulation motivating these finite-cutoff languages, in which the trivial Abelian squares \(aa,bb,cc\) are allowed.

Our basic question has a simple form:

\begin{quote}
\textbf{Among words that avoid all Abelian squares of half-length at most \(h\), how often is the entire word itself a longer Abelian square?}
\end{quote}

For \(XY\in L_h\cap\Sigma^{2n}\), \(|X|=|Y|=n\), set
\begin{equation}\label{eq:Rn}
R_n^{(h)}
=
\frac{
\#\{XY\in L_h\cap\Sigma^{2n}:\Psi(X)=\Psi(Y)\}
}{
|L_h\cap\Sigma^{2n}|
}.
\end{equation}
Thus \(R_n^{(h)}\) is literally the density of half-length-\(n\) Abelian squares inside the bounded-avoidance language \(L_h\).

The corresponding unrestricted problem is the natural external benchmark studied by Richmond and Shallit~\cite{Richmond2009}. Independently of that attribution, the free ternary multinomial model gives the leading constant
\[
C_{\rm free}
=
\frac{3\sqrt3}{4\pi}
=
0.413496\ldots.
\]
The exact normalization against the published Richmond--Shallit~\cite{Richmond2009} asymptotic is kept as a separate literature check rather than assumed here.
Our main result gives the constrained constants for the first five nontrivial cutoffs:

\begin{table}[h]
\centering
\begin{tabular}{rrr}
\toprule
\(h\) & \(C_h\) & \(C_h/C_{\rm free}\) \\
\midrule
2 & 0.616760 & 1.4916 \\
3 & 0.599007 & 1.4486 \\
4 & 0.569265 & 1.3767 \\
5 & 0.765537 & 1.8514 \\
6 & 0.996661 & 2.4103 \\
\bottomrule
\end{tabular}
\caption{The certified family of leading constants.}
\label{tab:family}
\end{table}

Hence every bounded-avoidance language in the certified family has a larger leading density of long Abelian squares than the free language, but the effect is not monotone:
\[
\boxed{C_2>C_3>C_4<C_5<C_6.}
\]

The first increase occurs at \(h=5\), which is also the first cutoff where the canonical window presentation loses recurrent states. The timing alone is only a finite-family structural fact. The stronger evidence comes from a \(K=5\) ablation: among the three newly rejected half-Parikh classes, type \((3,1,1)\) is the only one incident with the lost states, deleting it alone produces the final recurrent core, and its isolated collision constant is already essentially the full \(L_5\) value. The effects of the three classes are not additive.

The paper therefore has one main story:
\[
\text{non-monotone family}
\quad\longrightarrow\quad
\text{turn at }h=5
\quad\longrightarrow\quad
\text{anatomy of that turn}.
\]

The local-limit machinery itself is classical~\cite{Bertoni2003,Goldwurm2023,Herve2010}. Our contribution is the bounded-Abelian-avoidance application, the explicit family of leading constants, the non-monotone cutoff dependence, and the exact finite-state structure at the first upturn.

The work is motivated by the ternary Abelian-square avoidance question associated with M\"akel\"a; see Rao and Rosenfeld~\cite{Rao2018} for the closest positive avoidance result. Every fully avoiding infinite word would lie in every \(L_h\), but nothing proved here implies that such a word exists or settles the full conjecture.

\subsection*{Contributions}

\begin{enumerate}
\item \textbf{Family asymptotic.} For each \(h\in\{2,\ldots,6\}\),
\[
R_n^{(h)}\sim C_h/n
\]
with an explicit leading constant.
\item \textbf{Non-monotone cutoff dependence.} All five \(C_h\) exceed \(C_{\rm free}\), but the amplification weakens through \(h=4\) before increasing at \(h=5,6\).
\item \textbf{A quantitative \(K=5\) ablation and exact core-isolation theorem.} The \((3,1,1)\) transition class is tied to the first recurrent-state loss and the dominant change in \(C\); isolating the surviving core is an exact bipartite vertex-cover problem with minimum \(102\).
\end{enumerate}

\section{Bounded-avoidance languages and window presentations}
\label{sec:windows}

Set
\[
m=2h-1.
\]
A \textbf{valid state} is a ternary word of length \(m\) containing no Abelian square of half-length \(2,\ldots,h-1\). There is an edge
\[
u\longrightarrow v
\]
when the two windows overlap in \(m-1\) letters and appending the new letter does not create an Abelian square of half-length \(h\).

For every \(n\ge m\), words of \(L_h\cap\Sigma^n\) correspond bijectively to a choice of initial valid state together with a path of \(n-m\) edges in this graph. If \(A_h^{\rm full}\) denotes the full valid-state adjacency matrix, then
\begin{equation}\label{eq:pathword}
\boxed{
|L_h\cap\Sigma^n|
=
\mathbf1^T(A_h^{\rm full})^{\,n-m}\mathbf1.
}
\end{equation}

Starting from the valid graph, iteratively delete any state whose current in-degree or out-degree is zero. A state that survives is called \textbf{essential}. Equivalently, in a finite graph it can occur on a bi-infinite admissible path.

The raw, valid, and essential counts are:

\begin{table}[h]
\centering
\begin{tabular}{rrrrrr}
\toprule
\(h\) & \(m\) & raw \(3^m\) & valid & essential & essential edges \\
\midrule
2 & 3 & 27 & 27 & 27 & 66 \\
3 & 5 & 243 & 162 & 162 & 360 \\
4 & 7 & 2187 & 786 & 786 & 1572 \\
5 & 9 & 19683 & 3114 & 2844 & 5418 \\
6 & 11 & 177147 & 11070 & 10128 & 18774 \\
\bottomrule
\end{tabular}
\caption{Raw, valid and essential state counts of the memory-\(m\) window presentation.}
\label{tab:counts}
\end{table}

\begin{proposition}[Certified graph hypotheses]
\label{prop:graph}
For every \(h\in\{2,\ldots,6\}\):
\begin{enumerate}
\item the essential graph has a unique dominant strongly connected component;
\item that component has graph period \(1\);
\item the projected equal-length cycle Parikh-difference lattice is all of \(\mathbb Z^2\);
\item the associated covariance is nondegenerate;
\item the synchronized product of the time-reversed and forward core chains is primitive.
\end{enumerate}
\end{proposition}

The period and lattice claims are certified by explicit closed walks. In particular, for each \(h\), three equal-length closed walks at a common root have projected Parikh differences containing a unimodular pair. Item 5 follows because the forward and time-reversed core chains are primitive; the product of two primitive finite transition matrices is primitive as well.

The explicit roots, return lengths, cycles, Parikh vectors, and determinants are recorded in Appendix~A.

\begin{remark}[Why the non-essential part is acyclic]
\label{rem:acyclic}
Every vertex on a directed cycle has in-degree and out-degree at least \(1\) along that cycle at every stage of the pruning. Such a vertex can therefore never be deleted. Hence every directed cycle lies in the essential set, and the non-essential part of the valid graph is necessarily a DAG.

For \(h=2,3,4\) there is no non-essential part. For the full \(L_5\) presentation its longest directed path has \(D_5=3\) edges; for \(L_6\), \(D_6=5\).
\end{remark}

\section{Parry measure and the scalar variance \texorpdfstring{\(a_h\)}{a\_h}}
\label{sec:parry}

Let \(A_h\) be the adjacency matrix of the unique dominant essential component, with Perron root \(\lambda_h\) and positive right and left Perron vectors \(r_h,\ell_h\). The measure of maximal entropy is the Parry chain
\[
P_h(i,j)
=
\frac{(A_h)_{ij}r_h(j)}{\lambda_h r_h(i)}.
\]

The \(S_3\)-symmetry of the language forces stationary letter frequencies
\[
(1/3,1/3,1/3).
\]

Project the centered Parikh increment to two coordinates. Let \(\Sigma_h\) be its asymptotic covariance.

\begin{lemma}[\(S_3\)-equivariant covariance]
\label{lem:s3cov}
For every \(h\in\{2,\ldots,6\}\),
\[
\Sigma_h
=
a_h
\begin{pmatrix}
1&-\frac12\\
-\frac12&1
\end{pmatrix}
\]
for some \(a_h>0\).
\end{lemma}

\begin{proof}
The zero-sum Parikh subspace is the irreducible two-dimensional standard representation of \(S_3\). The asymptotic covariance commutes with the \(S_3\)-action and is therefore scalar on that representation. Writing the scalar operator in the coordinates \((a,b)\) gives the stated matrix.
\end{proof}

Consequently,
\[
\det\Sigma_h=\frac34a_h^2.
\]

The values used in this paper are:

\begin{table}[h]
\centering
\begin{tabular}{rrr}
\toprule
\(h\) & \(\lambda_h\) & \(a_h\) \\
\midrule
2 & 2.4511095375 & 0.1489852192 \\
3 & 2.2288029013 & 0.1534007516 \\
4 & 2.0666349657 & 0.1614154912 \\
5 & 1.9441605457 & 0.1200310311 \\
6 & 1.8483339782 & 0.0921960033 \\
\bottomrule
\end{tabular}
\caption{Perron roots and scalar Parikh variance rates.}
\label{tab:ah}
\end{table}

These values were obtained by a Green--Kubo/Poisson computation and independently checked by exact-moment dynamic programming. The displayed precision is intentionally shorter than the raw numerical output.

\section{Folding a long Abelian square}
\label{sec:fold}

Let
\[
w=w_0w_1\cdots w_{2n-1}\in L_h,
\qquad
X=w_0\cdots w_{n-1},
\qquad
Y=w_n\cdots w_{2n-1},
\]
and let
\[
Z_t=w_t\cdots w_{t+m-1}
\]
be the memory-\(m\) window.

Write
\[
N=n-m.
\]
The fold starts at the diagonal product state
\[
(Z_N,Z_N).
\]

Let \(\eta(u\to v)\in\mathbb Z^2\) be the projected Parikh vector of the new letter emitted by the forward edge \(u\to v\). For \(j=0,\ldots,N-1\), pair
\begin{itemize}
\item the left edge \(Z_{N-1-j}\to Z_{N-j}\), traversed backwards, and
\item the right edge \(Z_{N+j}\to Z_{N+j+1}\), traversed forwards.
\end{itemize}

Define the product increment
\[
g_j
=
\eta(Z_{N+j}\to Z_{N+j+1})
-
\eta(Z_{N-1-j}\to Z_{N-j}).
\]

\begin{lemma}[Exact fold identity]
\label{lem:fold}
With
\[
\beta
=
\Psi(Z_{2n-m})-\Psi(Z_0)
\]
projected to the same two coordinates,
\[
\boxed{
\Psi(Y)-\Psi(X)
=
\sum_{j=0}^{N-1}g_j+\beta.
}
\]
Moreover,
\[
\|\beta\|_\infty\le m.
\]
\end{lemma}

\begin{proof}
The paired edges account for
\[
w_m,\ldots,w_{n-1}
\]
on the left and
\[
w_n,\ldots,w_{2n-m-1}
\]
on the right. The only unpaired letters are the first \(m\) letters and the final \(m\) letters, whose Parikh difference is exactly
\[
\Psi(Z_{2n-m})-\Psi(Z_0).
\]
Summing the paired right-minus-left emissions and the boundary term gives \(\Psi(Y)-\Psi(X)\). Each state has length \(m\), giving the uniform bound.
\end{proof}

\begin{lemma}[Covariance doubling]
\label{lem:cov2}
Let \(\Sigma_{D,h}\) be the asymptotic covariance of the folded difference observable. Then
\[
\boxed{\Sigma_{D,h}=2\Sigma_h.}
\]
\end{lemma}

\begin{proof}
Condition on the center state \(Z_N=z\). Let \(S_N^+\) be the centered projected Parikh sum to the right of the center and \(S_N^-\) the corresponding sum on the left, read backwards. Given \(z\), the two are independent by the Markov property. The backward chain is the time reversal of the Parry chain, so it has the same asymptotic covariance:
\[
\operatorname{Cov}(S_N^\pm\mid z)
=
N\Sigma_h+O(1)
\]
uniformly over the finitely many center states.

Because the observable is centered at the stationary mean and the finite primitive chain forgets its initial state geometrically fast,
\[
\mathbb E[S_N^\pm\mid z]=O(1)
\]
uniformly in \(z\). The law of total covariance therefore gives
\[
\operatorname{Cov}(S_N^+-S_N^-)
=
2N\Sigma_h+O(1).
\]
Divide by \(N\) and let \(N\to\infty\).
\end{proof}

The same matrix \(\Sigma_{D,h}\) is the Hessian covariance appearing in the logarithm of the leading eigenvalue of the Fourier-tilted product transfer matrix introduced below; this is the standard Markov-additive covariance/pressure identification~\cite{Herve2010}.

\section{Endpoint-uniform LLT for the product chain}
\label{sec:llt}

The fold evolves on the product of the time-reversed left chain and the forward right chain. Let
\[
M_h(\theta)
\]
denote the Fourier-tilted product transfer matrix for the rank-two increment \(g\).

The analytic input needed here is an endpoint-uniform lattice local limit theorem. Uniformly over product-chain endpoints, bounded target displacements \(v\in\mathbb Z^2\), and lengths \(N+O(1)\),
\begin{equation}\label{eq:llt}
\mathcal N_N(\alpha,\omega;v)
=
\mathcal W_N(\alpha,\omega)
\frac{1}{2\pi N\sqrt{\det\Sigma_{D,h}}}
(1+o(1)),
\end{equation}
where \(\mathcal W_N\) is the unconstrained product-path weight/count for those endpoints. The full finite-state proof is given in Appendix~F.

\subsection*{Product-lattice certificate}

The lattice hypothesis for the product observable follows directly from Proposition~\ref{prop:graph}. Let
\[
\ell_1,\ell_2,\ell_3
\]
be equal-length closed walks at a common single-chain root \(z_0\), with two unimodular projected Parikh differences. Pair the reversal of \(\ell_i\) on the left with \(\ell_j\) on the right. This is a product-chain cycle at \((z_0,z_0)\) whose \(g\)-sum is
\[
\Psi(\ell_j)-\Psi(\ell_i).
\]
Choosing \(i=j\) gives \(0\), while appropriate choices of \(i,j\) give the same unimodular pair certified in Proposition~\ref{prop:graph}. Hence the product observable has full lattice span \(\mathbb Z^2\).

\subsection*{Spectral proof scheme}

For completeness, the finite-state proof uses:
\begin{enumerate}
\item Fourier inversion on \(\mathbb Z^2\);
\item spectral decay of \(M_h(\theta)\) away from \(\theta=0\);
\item analytic perturbation of the simple leading eigenvalue near \(0\);
\item zero drift from the \(1/3\) stationary frequencies;
\item quadratic term \(\Sigma_{D,h}\);
\item the product-lattice certificate above to rule out nontrivial peripheral characters.
\end{enumerate}

This yields the displayed LLT uniformly in the finitely many endpoints and bounded displacements needed after transient peeling.

These are the standard finite-state Fourier/spectral steps; the project-specific work is the verification of primitivity, zero drift, nondegenerate covariance, full product lattice, and the endpoint-uniform bounded-displacement regime. The manuscript does not rely on a novelty claim for the LLT itself.

\section{Transient cancellation and the main theorem}
\label{sec:cancellation}

For this section only, write
\begin{itemize}
\item \(E_h\) for the dominant essential component of the full \(L_h\) presentation, and
\item \(T_h\) for the non-essential states of that same \(L_h\) graph.
\end{itemize}

This notation is \textbf{not} the CORE/LOST partition used later for the \(L_4\) scissor problem.

By Remark~\ref{rem:acyclic}, \(T_h\) is a DAG. A path in the full \(L_h\) graph can meet \(T_h\) only in a bounded prefix or suffix. Indeed, a genuine
\[
E_h\to T_h\to E_h
\]
excursion would place the traversed \(T_h\)-vertices on a directed cycle, contradicting non-essentiality.

\begin{lemma}[Exact full-path identity]
\label{lem:pathid}
For every \(n\ge m\), with \(N=n-m\),
\[
\boxed{
\sum_{z,s,u,e}
\bigl((A_h^{\rm full})^N\bigr)_{s,z}
\bigl((A_h^{\rm full})^N\bigr)_{z,u}
\bigl((A_h^{\rm full})^m\bigr)_{u,e}
=
\mathbf1^T(A_h^{\rm full})^{2n-m}\mathbf1
=
|L_h\cap\Sigma^{2n}|.
}
\]
\end{lemma}

\begin{proof}
This is pure path algebra. Splitting a length-\((2N+m)\) full-graph path after \(N\) and \(2N\) edges is a bijection with the triples counted on the left. Since \(2N+m=2n-m\), the path-word bijection~\eqref{eq:pathword} of Section~\ref{sec:windows} gives the last equality.
\end{proof}

\begin{lemma}[Leading boundary factor]
\label{lem:boundary}
For \(h\in\{2,\ldots,6\}\), the leading all-language boundary multiplier in the folded LLT is
\[
\boxed{B_h=1.}
\]
\end{lemma}

\begin{proof}
For \(h=2,3,4\), \(T_h=\varnothing\), so no transient peeling is needed.

For \(h=5,6\), let \(D_h\) be the maximum path length in \(T_h\), with
\[
D_5=3,\qquad D_6=5.
\]
Consider one term in the exact decomposition of Lemma~\ref{lem:pathid}. By Lemma~\ref{lem:fold}, its endpoint pair determines one required displacement,
\[
v=\Psi(Z_0)-\Psi(Z_{2n-m}),
\qquad
\|v\|_\infty\le m.
\]
Peel any transient prefix and suffix. At most \(D_h\) edges are removed at either end, so the remaining synchronized product path lies in the core, has length
\[
N-O(1),
\]
and its required displacement differs from \(v\) by only \(O(1)\). Thus the bounded-displacement hypothesis of Section~\ref{sec:llt} applies uniformly.

Apply the endpoint-uniform product-chain LLT~\eqref{eq:llt} to this peeled core term. Uniformity in endpoints, in bounded displacement, and in lengths \(N+O(1)\) gives the same leading density
\[
D_N
=
\frac{1}{2\pi N\sqrt{\det\Sigma_{D,h}}}
\]
with a relative error \(o(1)\) uniform over all endpoint configurations.

Thus each nonnegative term in the exact full-path decomposition equals its unconstrained path multiplicity times
\[
D_N(1+o(1)),
\]
with one uniform relative error. Summing preserves the relative \(o(1)\). Lemma~\ref{lem:pathid} then gives
\[
\text{numerator}
=
D_N\,|L_h\cap\Sigma^{2n}|(1+o(1)).
\]
Therefore \(B_h=1\).
\end{proof}

The uniqueness of the dominant SCC is used here to obtain one common asymptotic density. With two dominant components carrying the same Perron root, different covariance data could prevent a single factor from being extracted uniformly.

\begin{theorem}[Abelian-square density in the certified family]
\label{thm:main}
For every
\[
h\in\{2,3,4,5,6\},
\]
\[
\boxed{
R_n^{(h)}
\sim
\frac{C_h}{n},
\qquad
C_h
=
\frac{1}{4\pi\sqrt{\det\Sigma_h}}
=
\frac{1}{2\sqrt3\,\pi\,a_h}.
}
\]
\end{theorem}

\begin{proof}
By Proposition~\ref{prop:graph} the synchronized product process is primitive, and the product-lattice certificate of Section~\ref{sec:llt} gives full lattice span. Lemma~\ref{lem:fold} folds the equal-Parikh event into a bounded-displacement rank-two additive event. Lemma~\ref{lem:cov2} gives
\[
\Sigma_{D,h}=2\Sigma_h.
\]
The endpoint-uniform product LLT at synchronized length
\[
N=n-m
\]
therefore contributes
\[
\frac{1}{2\pi N\sqrt{\det(2\Sigma_h)}}.
\]
Lemma~\ref{lem:boundary} gives boundary factor \(1\). Since
\[
\frac1N=\frac1n(1+o(1)),
\]
we obtain
\[
R_n^{(h)}
\sim
\frac{1}{4\pi n\sqrt{\det\Sigma_h}}.
\]
Finally \(\det\Sigma_h=3a_h^2/4\).
\end{proof}

\begin{corollary}[Free-language comparison]
\label{cor:free}
For the unrestricted ternary language,
\[
a_{\rm free}=\frac29,
\qquad
C_{\rm free}=\frac{3\sqrt3}{4\pi}.
\]
Hence
\[
\boxed{
\frac{C_h}{C_{\rm free}}
=
\frac{2/9}{a_h}
}
\]
for every \(h\in\{2,\ldots,6\}\). In particular,
\[
C_h>C_{\rm free}
\]
throughout the certified family.
\end{corollary}

\begin{corollary}[Non-monotonicity]
\label{cor:nonmono}
\[
\boxed{
C_2>C_3>C_4<C_5<C_6.
}
\]
\end{corollary}

The first increase occurs at \(h=5\), the same cutoff at which valid and essential state counts first differ. No general causal law is claimed.

\begin{remark}[Unrestricted benchmark]
\label{rem:benchmark}
The free-language constant \(C_{\rm free}=3\sqrt3/(4\pi)\) is obtained directly from the unrestricted ternary multinomial model and is reproduced by the same covariance formula with \(a_{\rm free}=2/9\). Richmond and Shallit's~\cite{Richmond2009} unrestricted Abelian-square enumeration is therefore the natural published benchmark. Matching the precise normalization of their stated asymptotic to the present \(R_n\)-notation is a literature-verification step, not an assumption used in Theorem~\ref{thm:main}.
\end{remark}

\section{Anatomy of the \texorpdfstring{\(h=4\to5\)}{h=4 to 5} turn}
\label{sec:anatomy}

We now switch to a \textbf{different graph object}: the common-memory-\(9\) presentation of \(L_4\), used to ask which newly rejected \(K=5\) transitions isolate the eventual \(L_5\) core.

To prevent confusion with Section~\ref{sec:cancellation}, write
\begin{itemize}
\item \(\mathrm{CORE}\) for the \(2844\)-state set that survives the full \(K=5\) restriction;
\item \(\mathrm{LOST}=V(L_4)\setminus\mathrm{CORE}\) for the \(270\)-state complement.
\end{itemize}

In the pre-deletion \(L_4\) graph, every LOST state is both reachable from CORE and able to return to CORE. Thus CORE\(\to\)LOST\(\to\)CORE excursions are exactly the phenomenon studied here. This is the opposite of the transient geometry of \(T_h\) in the already-essentialized \(L_h\) graph of Section~\ref{sec:cancellation}.

The newly rejected \(K=5\) transitions split by half-Parikh type:

\begin{table}[h]
\centering
\begin{tabular}{lrrrrr}
\toprule
type & total & CORE\(\to\)CORE & CORE\(\to\)LOST & LOST\(\to\)CORE & LOST\(\to\)LOST \\
\midrule
\((2,2,1)\) & 222 & 222 & 0 & 0 & 0 \\
\((3,1,1)\) & 216 & 72 & 66 & 66 & 12 \\
\((3,2,0)\) & 12 & 12 & 0 & 0 & 0 \\
\bottomrule
\end{tabular}
\caption{Census of the \(K=5\) rejected transitions by half-Parikh type.}
\label{tab:census}
\end{table}

Only \((3,1,1)\) is incident with LOST.

\subsection*{\texorpdfstring{\(K=5\)}{K=5} ablations}

Let \(G_{320}\), \(G_{311}\), and \(G_{221}\) denote the common-memory graphs obtained by deleting only the indicated half-Parikh class.

\begin{table}[h]
\centering
\begin{tabular}{lrrr}
\toprule
graph & essential & edges & \(C\) \\
\midrule
\(L_4\) & 3114 & 6300 & 0.569265 \\
\(G_{320}\) & 3114 & 6288 & 0.587243 \\
\(G_{311}\) & 2844 & 5652 & 0.768856 \\
\(G_{221}\) & 3114 & 6078 & 0.514383 \\
\(L_5\) & 2844 & 5418 & 0.765537 \\
\bottomrule
\end{tabular}
\caption{Single-class \(K=5\) ablations in the common memory-\(9\) presentation.}
\label{tab:ablation}
\end{table}

Deleting only \((3,1,1)\) produces the same essential state set as the full \(K=5\) restriction.

Define
\[
\Delta_{320}=C(G_{320})-C(L_4)=+0.0179785\ldots,
\]
\[
\Delta_{311}=C(G_{311})-C(L_4)=+0.1995914\ldots,
\]
and
\[
\Delta_{221}=C(G_{221})-C(L_4)=-0.0548822\ldots.
\]
The full change is
\[
\Delta_{\rm full}
=
C(L_5)-C(L_4)
=
0.1962719\ldots.
\]

The single-class changes sum to
\[
0.1626876\ldots,
\]
so the interaction term
\[
\boxed{
\Delta_{\rm int}
=
\Delta_{\rm full}
-
(\Delta_{320}+\Delta_{311}+\Delta_{221})
=
+0.0335843\ldots
}
\]
is nonzero. In particular, the \((2,2,1)\) class alone moves the constant in the opposite direction from the dominant \((3,1,1)\) change.

Also,
\[
\frac{\Delta_{311}}{\Delta_{\rm full}}
=
1.016913\ldots.
\]
This ratio is descriptive, not an additive attribution of cause. The safe conclusion is that the \((3,1,1)\) ablation is associated with both the entire recurrent-state loss and a collision-constant change slightly larger than the full \(K=5\) change, while the remaining deletions and their interactions partially offset it.

The induced LOST graph in the pre-deletion \(L_4\) graph is acyclic, with \(216\) internal edges and longest directed path \(5\) edges. After the full \(K=5\) deletion, the non-essential \(L_5\) subgraph instead has \(204\) internal edges and longest path \(3\). These are different graphs.

\section{Excursion cover and the exact scissor}
\label{sec:scissor}

Let
\[
E_{\rm in}
=
E\cap(\mathrm{CORE}\times\mathrm{LOST}),
\qquad
E_{\rm out}
=
E\cap(\mathrm{LOST}\times\mathrm{CORE}).
\]
The full boundary has
\[
|E_{\rm in}|=|E_{\rm out}|=180.
\]

Construct a bipartite conflict graph \(H\) whose left vertices are \(E_{\rm in}\) and right vertices are \(E_{\rm out}\). Join \(e\in E_{\rm in}\) to \(f\in E_{\rm out}\) whenever
\[
\operatorname{head}(e)
\leadsto
\operatorname{tail}(f)
\]
inside the induced LOST graph, allowing the empty path.

A vertex of \(H\) therefore represents a boundary \textbf{transition} of the original automaton.

\begin{proposition}[Excursion-cover equivalence]
\label{prop:excursion}
Let \(G\) be a finite directed graph, let \(C^\star\) induce a strongly connected subgraph containing a directed cycle, and let \(T^\star=V(G)\setminus C^\star\) induce an acyclic subgraph. Let \(E_{\rm in}\) and \(E_{\rm out}\) be \textbf{all} boundary edges between \(C^\star\) and \(T^\star\), and form \(H\) as above.

For
\[
S\subseteq E_{\rm in}\cup E_{\rm out},
\]
\[
\boxed{
\mathrm{Ess}(G-S)=C^\star
\iff
S\text{ is a vertex cover of }H.
}
\]
\end{proposition}

\begin{proof}
Suppose \(S\) covers \(H\). If a vertex of \(T^\star\) were both reachable from \(C^\star\) and able to return to \(C^\star\) in \(G-S\), there would be a surviving entry edge \(e\), a surviving exit edge \(f\), and an internal \(T^\star\)-path connecting them. Then \(e\) and \(f\) would form an uncovered edge of \(H\), contradiction. Since \(G[T^\star]\) is acyclic, no vertex of \(T^\star\) can remain essential. The subgraph induced by \(C^\star\) is untouched because only boundary edges are deleted.

Conversely, if \(S\) is not a vertex cover, some adjacent entry \(e\) and exit \(f\) both survive. Their adjacency gives an internal \(T^\star\)-path. Strong connectivity of \(C^\star\) closes this excursion through \(C^\star\) into a recurrent directed walk, making every vertex on the \(T^\star\)-segment essential. Hence the essential set is strictly larger than \(C^\star\).
\end{proof}

By K\H{o}nig's theorem,
\[
\tau(H)=\nu(H).
\]

For the \(K=5\) scissor graph,
\[
|V_L(H)|=|V_R(H)|=180,
\qquad
|E(H)|=366,
\]
and explicit matching and cover certificates (Appendix~C) give
\[
\boxed{\tau(H)=\nu(H)=102.}
\]

\subsection*{Restricted \texorpdfstring{\((3,1,1)\)}{(3,1,1)} deletion pool}

Only
\[
66+66=132
\]
of the \(360\) full CORE-LOST boundary transitions belong to the \(K=5\)-rejected \((3,1,1)\) pool. Restricting deletions to this pool is therefore a \textbf{vertex-cover problem with forbidden vertices}, not literally the unrestricted setting of Proposition~\ref{prop:excursion}.

Nevertheless the restricted optimum remains \(102\). A human-readable certificate is
\[
\boxed{84+18=102}.
\]
Here \(84\) pool transitions are forced because their conflict partner lies outside the deletable pool. After removing them, the residual conflict graph contains \(18\) disjoint copies of \(K_2\), each requiring exactly one of its two conflict-graph vertices. Since each conflict-graph vertex represents a boundary transition, these contribute exactly \(18\) further deletions.

Twelve remaining pool transitions are redundant for minimum covers: they are not isolated in \(H\), but every excursion incident with them is already broken by forced deletions.

\subsection*{Symmetry refinements}

Under \(S_3\) invariance the optimum is \(17\) transition-orbits, i.e.\ \(102\) transitions. Under \(S_3\) plus reversal invariance the minimum is \(9\) reversal pairs, i.e.\ \(108\) transitions. Detailed orbit enumeration and counts of optimal solutions are given in Appendix~D.

\section{Numerical validation and reproducibility}
\label{sec:numerics}

Finite-\(n\) dynamic programming is used only as an independent check of the theorem-derived constants.

For \(h=2,3,4\), independent all-language finite-\(n\) computations and second-order extrapolations agree with the spectral constants to approximately \(0.003\%\), \(0.009\%\), and \(0.012\%\), respectively.

For \(h=5\), the corrected all-language values include:

\begin{table}[h]
\centering
\begin{tabular}{rr}
\toprule
\(n\) & \(nR_n^{(5)}\) \\
\midrule
20 & 0.6066545 \\
30 & 0.6583281 \\
40 & 0.6833348 \\
50 & 0.6992042 \\
60 & 0.7099965 \\
\bottomrule
\end{tabular}
\caption{All-language finite-\(n\) control for \(h=5\).}
\label{tab:h5}
\end{table}

\noindent with theorem constant
\[
C_5=0.7655366150\ldots.
\]
At \(n=60\) the finite-\(n\) value \(0.7099965\) is approximately \(7.3\%\) below \(C_5\); a second-order Richardson extrapolation of the same series lies approximately \(0.13\%\) from \(C_5\). Finite-\(n\) data validate, but do not prove, the asymptotic.

A previous series was accidentally core-only; the values above are the corrected all-language control.

\subsection*{\texorpdfstring{\(h=6\)}{h=6} finite-\texorpdfstring{\(n\)}{n} validation}

The theorem-derived value
\[
C_6=0.9966608746\ldots
\]
rests on the independently certified graph, lattice, covariance, and local-limit derivation. Finite-\(n\) data are used only as an independent numerical check. The full evidence-level breakdown for \(h=6\) is given in Appendix~B.

\begin{table}[h]
\centering
\begin{tabular}{rrrr}
\toprule
\(n\) & \(|L_6\cap\Sigma^{2n}|\) & collisions & \(nR_n^{(6)}\) \\
\midrule
12 & 32,264,706 & 1,482,864 & 0.5515118594 \\
20 & 598,401,288,060 & 19,842,060,168 & 0.6631690327 \\
30 & 129,601,490,959,197,930 & 3,289,168,271,775,816 & 0.7613727853 \\
40 & 28,068,484,992,703,808,552,340 & 570,703,655,838,606,639,024 & 0.8133016883 \\
50 & 6,078,940,846,562,758,147,891,532,058 & 102,845,141,999,403,991,617,566,730 & 0.8459133309 \\
\bottomrule
\end{tabular}
\caption{Verified raw all-language finite-\(n\) values for \(h=6\).}
\label{tab:h6}
\end{table}

\subsection*{Reproducibility}

Reproducibility materials comprise the graph builders, SCC and graph-period certificates, explicit single- and product-lattice witnesses, two independent covariance calculations, all-language finite-\(n\) controls, \(K=5\) ablation and cover certificates, symmetry calculations, and source hashes/commands/environments. The preserved \(h=6\) reporting defect and its arithmetic correction audit are included with these materials. Appendix~E records the provenance level of each item.

\section{Discussion, relation to M\"akel\"a, and open questions}
\label{sec:discussion}

The family gives two simultaneous facts:
\[
C_h>C_{\rm free}
\qquad(h=2,\ldots,6),
\]
yet
\[
C_2>C_3>C_4<C_5<C_6.
\]
Thus bounded local avoidance makes long Abelian squares more common than in the free language throughout the certified family, but increasing the cutoff does not increase that effect monotonically.

The identity
\[
\frac{C_h}{C_{\rm free}}
=
\frac{2/9}{a_h}
\]
shows that the phenomenon is exactly controlled by the scalar asymptotic Parikh variance.

At \(h=5\), the first upturn occurs at the same cutoff as the first valid-to-essential state loss. The ablation goes further than this timing coincidence: the \((3,1,1)\) class is the only new rejection class incident with LOST, deleting it alone produces the final core, and its isolated \(C\)-shift dominates the observed change. Because the ablations are non-additive, we do not claim a general causal law.

The Perron roots decrease with \(h\),
\[
2.4511,\ 2.2288,\ 2.0666,\ 1.9442,\ 1.8483,
\]
while the collision constants do not. Entropy reduction and concentration of the Parikh distribution therefore need not move in lockstep.

\subsection*{Relation to M\"akel\"a's problem}

Every infinite ternary word avoiding all nontrivial Abelian squares would lie in every finite cutoff language \(L_h\). Rao and Rosenfeld~\cite{Rao2018} established the closest positive avoidance result, for sufficiently large periods. The present results describe necessary finite-window containers, but they do not prove:
\begin{itemize}
\item existence of an infinite word surviving every cutoff;
\item convergence of the finite-cutoff Parry measures;
\item a positive growth rate for the full infinite-avoidance language;
\item a proof or disproof of M\"akel\"a's conjecture.
\end{itemize}

\subsection*{Related work and novelty boundary}

Richmond and Shallit~\cite{Richmond2009} provide the direct unrestricted Abelian-square enumeration benchmark. General symbol-count central and local limit theorems for regular languages are treated by Bertoni et al.~\cite{Bertoni2003} and later work~\cite{Goldwurm2023}; multidimensional spectral Markov methods are standard~\cite{Herve2010}. Directional counting for regular languages is developed by Grigorchuk and Quint~\cite{Grigorchuk2026}. Repetition-free language literature~\cite{Shur2012,Petrova2021} supplies the surrounding growth and avoidance context.

The paper does \textbf{not} claim a new general local-limit theorem. At the present audit stage the strongest safe novelty wording remains:

\begin{quote}
We are not aware of a previous asymptotic enumeration of adjacent Abelian-square collisions inside a finite-type language defined by bounded Abelian-square avoidance.
\end{quote}

Before submission, the citation audit should verify the exact normalization of the published Richmond--Shallit~\cite{Richmond2009} asymptotic against the present \(R_n\)-notation and complete the final citation-neighborhood search.

\subsection*{Open questions}

The present paper leaves outside its theorem package:
\begin{enumerate}
\item the behavior of \(C_h\) for \(h\ge7\);
\item whether later recurrent-core losses accompany sharp changes in \(a_h\);
\item whether analogous scissor structures recur at later cutoffs;
\item second-order terms in \(R_n^{(h)}\);
\item larger alphabets and non-adjacent blocks.
\end{enumerate}

No \(h=7\) computation is needed before the present manuscript is stabilized.

\bibliographystyle{plain}
\bibliography{references}

\appendix
\section{Period and lattice certificates}
\label{app:A}

This appendix records the explicit witnesses for parts~2 and~3 of
Proposition~\ref{prop:graph}. All strings below are the sequences of
\emph{emitted letters} of a closed walk read from the stated root; the walk
returns to that root. Every walk was checked edge by edge against the
memory-\(m\) window graph, and the state validity of the root was checked
independently.

For each \(h\) the period certificate consists of two closed walks at the root
whose lengths are coprime, which forces graph period \(1\). The lattice
certificate consists of three closed walks of one common length \(L\) at the
same root; writing \(\Psi\) for the Parikh vector and projecting onto the first
two coordinates, the two differences
\(d_1=\Psi(c_2)-\Psi(c_1)\) and \(d_2=\Psi(c_3)-\Psi(c_1)\) satisfy
\(|\det[d_1\ d_2]|=1\), so they generate \(\mathbb Z^2\) and the projected
lattice has covolume \(1\).

\subsection*{\texorpdfstring{\(h=2\)}{h=2}, \texorpdfstring{\(m=3\)}{m=3}, root \texttt{aaa}}
\begin{itemize}
\item Period: closed walks of lengths \(4\) (\texttt{baaa}) and \(5\)
      (\texttt{bcaaa}); \(\gcd(4,5)=1\).
\item Lattice at \(L=6\):
      \(c_1=\texttt{bacaaa}\), \(\Psi=(4,1,1)\);
      \(c_2=\texttt{bbbaaa}\), \(\Psi=(3,3,0)\);
      \(c_3=\texttt{bbcaaa}\), \(\Psi=(3,2,1)\).
\item \(d_1=(-1,2)\), \(d_2=(-1,1)\), \(\det=1\).
\end{itemize}

\subsection*{\texorpdfstring{\(h=3\)}{h=3}, \texorpdfstring{\(m=5\)}{m=5}, root \texttt{aabaa}}
\begin{itemize}
\item Period: closed walks of lengths \(6\) (\texttt{caabaa}) and \(7\)
      (\texttt{acaabaa}); \(\gcd(6,7)=1\).
\item Lattice at \(L=9\):
      \(c_1=\texttt{acbcaabaa}\), \(\Psi=(5,2,2)\);
      \(c_2=\texttt{acccaabaa}\), \(\Psi=(5,1,3)\);
      \(c_3=\texttt{cbccaabaa}\), \(\Psi=(4,2,3)\).
\item \(d_1=(0,-1)\), \(d_2=(-1,0)\), \(\det=-1\).
\end{itemize}

\subsection*{\texorpdfstring{\(h=4\)}{h=4}, \texorpdfstring{\(m=7\)}{m=7}, root \texttt{aaabaaa}}
\begin{itemize}
\item Period: closed walks of lengths \(8\) (\texttt{caaabaaa}) and \(11\)
      (\texttt{cabcaaabaaa}); \(\gcd(8,11)=1\).
\item Lattice at \(L=13\):
      \(c_1=\texttt{caaabcaaabaaa}\), \(\Psi=(9,2,2)\);
      \(c_2=\texttt{caabccaaabaaa}\), \(\Psi=(8,2,3)\);
      \(c_3=\texttt{cabbccaaabaaa}\), \(\Psi=(7,3,3)\).
\item \(d_1=(-1,0)\), \(d_2=(-2,1)\), \(\det=-1\).
\end{itemize}

\subsection*{\texorpdfstring{\(h=5\)}{h=5}, \texorpdfstring{\(m=9\)}{m=9}, root \texttt{aaabaaaca}}
\begin{itemize}
\item Period: closed walks of lengths \(8\) (\texttt{aabaaaca}) and \(13\)
      (\texttt{abccaaabaaaca}); \(\gcd(8,13)=1\).
\item Lattice at \(L=13\):
      \(c_1=\texttt{abccaaabaaaca}\), \(\Psi=(8,2,3)\);
      \(c_2=\texttt{bbccaaabaaaca}\), \(\Psi=(7,3,3)\);
      \(c_3=\texttt{bcccaaabaaaca}\), \(\Psi=(7,2,4)\).
\item \(d_1=(-1,1)\), \(d_2=(-1,0)\), \(\det=1\).
\end{itemize}

\subsection*{\texorpdfstring{\(h=6\)}{h=6}, \texorpdfstring{\(m=11\)}{m=11}, root \texttt{abbbcaaaccb}}
\begin{itemize}
\item Period: closed walks of lengths \(12\) (\texttt{babbbcaaaccb}) and \(17\)
      (\texttt{acccaaabbbcaaaccb}); \(\gcd(12,17)=1\).
\item Lattice at \(L=17\):
      \(c_1=\texttt{acccbbabbbcaaaccb}\), \(\Psi=(5,6,6)\);
      \(c_2=\texttt{baaacaabbbcaaaccb}\), \(\Psi=(8,5,4)\);
      \(c_3=\texttt{bbccaaabbbcaaaccb}\), \(\Psi=(6,6,5)\).
\item \(d_1=(3,-1)\), \(d_2=(1,0)\), \(\det=1\).
\end{itemize}

\subsection*{Root independence}

The projected lattice generated by differences of equal-length closed walks
does not depend on the root chosen. If \(c,c'\) are equal-length closed walks
at \(z\), and \(p\) is a path \(z'\to z\) and \(q\) a path \(z\to z'\) inside
the same strongly connected component, then \(pcq\) and \(pc'q\) are
equal-length closed walks at \(z'\) with the same Parikh difference. A single
root therefore suffices for each \(h\).

\section{Finite-\texorpdfstring{\(n\)}{n} validation and evidence levels}
\label{app:B}

Finite-\(n\) dynamic programming validates the theorem-derived constants; it
does not prove them. This appendix records exactly which quantities received
which kind of check.

\subsection*{\texorpdfstring{\(h=2,3,4\)}{h=2,3,4}}

Independent all-language finite-\(n\) computations, followed by second-order
Richardson extrapolation of \(nR_n^{(h)}\), agree with the spectral constants
to approximately \(0.003\%\) (\(h=2\)), \(0.009\%\) (\(h=3\)) and \(0.012\%\)
(\(h=4\)).

\subsection*{\texorpdfstring{\(h=5\)}{h=5}}

The all-language series is given in Table~\ref{tab:h5}. It approaches \(C_5\)
slowly: at \(n=60\) the value \(nR_n^{(5)}=0.7099965\) lies approximately
\(7.3\%\) below \(C_5=0.7655366150\ldots\). A second-order Richardson
extrapolation of the same series lies approximately \(0.13\%\) from \(C_5\).
We therefore do \emph{not} claim that the finite-\(n\) series reproduces the
constant directly; it is consistent with it after extrapolation.

An earlier series was accidentally computed on the essential core only. The
values in Table~\ref{tab:h5} are the corrected all-language control and must
not be conflated with the core-only sequence.

\subsection*{\texorpdfstring{\(h=6\)}{h=6}: the reporting-only defect and its evidence levels}

For \(h=6\), an earlier prose table reported absolute counts that did not
correspond to the preserved raw output of the collision dynamic program. A
subsequent arithmetic and provenance audit established that this was a
reporting-only defect, and the evidence separates into three distinct levels.
First, all five denominators \(|L_6\cap\Sigma^{2n}|\) for
\(n=12,20,30,40,50\) were independently reproduced by a separate exact
arbitrary-precision language-count dynamic program. Second, the collision
numerators for \(n=12,20,30\) were independently reproduced by a separate
arbitrary-precision implementation. Third, for \(n=40\) and \(n=50\) the
numerators were \textbf{not} recomputed by a second independent
implementation; instead the raw computation was audited for arithmetic storage
safety, the largest observed typed-array cell being \(2\,182\,364\,991\,314\),
far below the \(\mathtt{BigUint64Array}\) bound
\(2^{64}-1=18\,446\,744\,073\,709\,551\,615\), so no wraparound can have
occurred. All five reported values of \(nR_n^{(6)}\) are internally consistent
with the preserved raw counts as exact rationals, and the ratios reported
before the correction were unchanged by it. We make no claim that the
\(n=40,50\) numerators received a second independent computation.

A small-\(n\) brute-force control over \(\Sigma^{\le 13}\) and a regression of
the same pipeline against the previously validated \(h=5\) all-language series
both pass.

The verified raw values are reproduced in Table~\ref{tab:h6}. The correction
affects only the earlier prose report, not the theorem and not the raw
\(h=6\) validation data.

\section{Excursion-cover certificates}
\label{app:C}

All quantities in this appendix are produced and checked by the package-local
script \texttt{verify\_scissor\_and\_symmetry.js}, which rebuilds the memory-9
window graph from first principles and reads no external data file. Its
output for the present package is reproduced verbatim in
\texttt{VERIFIER\_OUTPUT.txt}; the explicit matching and cover are written to
\texttt{excursion\_cover\_certificate.json}.

\subsection*{Graph and boundary}

\begin{center}
\begin{tabular}{lrr}
\toprule
quantity & value & checked \\
\midrule
memory-9 valid states & 3114 & yes \\
\(L_4\) edges (all overlaps) & 6300 & yes \\
CORE & 2844 & yes \\
LOST & 270 & yes \\
LOST-internal edges in \(L_4\) & 216 & yes \\
\(|E_{\rm in}|\) & 180 & yes \\
\(|E_{\rm out}|\) & 180 & yes \\
total boundary & 360 & yes \\
\(|E(H)|\) & 366 & yes \\
\bottomrule
\end{tabular}
\end{center}

\subsection*{Optimum, with both certificates}

\[
\nu(H)=102,\qquad \tau(H)=102 .
\]

Two independent certificates are supplied and machine-checked:

\begin{itemize}
\item a \textbf{matching} of \(102\) pairwise vertex-disjoint edges of \(H\),
      each verified to be an actual edge of \(H\); this gives
      \(\tau(H)\ge\nu(H)\ge102\);
\item a \textbf{vertex cover} of \(102\) vertices (66 on the \(E_{\rm in}\)
      side, 36 on the \(E_{\rm out}\) side), verified to leave \(0\) of the
      \(366\) edges of \(H\) uncovered; this gives \(\tau(H)\le102\).
\end{itemize}

Together the two certificates pin the optimum exactly, independently of
K\H{o}nig's theorem. Each cover and matching vertex is stored in the JSON
certificate as the pair of 9-letter windows naming the boundary transition, so
the certificate can be re-checked against any independent construction of the
graph.

\subsection*{Restricted \texorpdfstring{\((3,1,1)\)}{(3,1,1)} pool}

\begin{center}
\begin{tabular}{lrr}
\toprule
quantity & value & checked \\
\midrule
pool size & 132 & yes \\
\(H\)-edges with no pool endpoint & 0 & yes \\
forced (conflict partner outside the pool) & 84 & yes \\
residual components & 18 & yes \\
all residual components are \(K_2\) & true & yes \\
redundant pool transitions & 12 & yes \\
restricted minimum \(=84+18\) & 102 & yes \\
\bottomrule
\end{tabular}
\end{center}

The zero in the second row is the feasibility check: every edge of \(H\) has at
least one endpoint inside the pool, so the forbidden-vertex problem of
Section~\ref{sec:scissor} is solvable at all. The decomposition
\(132=84+36+12\) accounts for every pool transition.

\section{Symmetry and reversal orbits}
\label{app:D}

The quantities below are produced and checked by the same package-local script
\texttt{verify\_scissor\_and\_symmetry.js}; the orbit representatives and the
reversal pairing are written to \texttt{symmetry\_orbits\_certificate.json}.
No external data file is read.

\subsection*{\texorpdfstring{\(S_3\)}{S3} action on the pool}

The alphabet group \(S_3\) acts on the memory-9 window graph by permuting
letters, and this action preserves the \((3,1,1)\) boundary pool. The 132 pool
transitions decompose into \textbf{22 orbits, every one of size 6}; both facts
are checked, and closure of the action on the pool is asserted by the script
(it raises an error otherwise).

Restricting deletions to unions of whole orbits gives a vertex-cover problem
on the 22-vertex quotient conflict graph. Exhaustive search over all
\(2^{22}\) orbit subsets yields

\[
\text{minimum }=17\text{ orbits}=17\times6=102\text{ transitions},
\qquad
\#\{\text{optima}\}=8 .
\]

The \(S_3\)-invariant optimum therefore attains the same value 102 as the
unrestricted optimum of Appendix~\ref{app:C}; symmetry costs nothing here.

\subsection*{Reversal}

Reversal is a directed \emph{anti}-automorphism: an edge \(u\to v\) is carried
to \(\operatorname{rev}(v)\to\operatorname{rev}(u)\). It therefore exchanges
\(E_{\rm in}\) and \(E_{\rm out}\), and no orbit can be fixed by it. The script
checks that the induced map on orbits is an involution and that the 22 orbits
form \textbf{11 reversal pairs}.

Requiring the deletion set to be invariant under \(S_3\) \emph{and} reversal
means selecting whole reversal pairs. Exhaustive search over all \(2^{11}\)
selections yields

\[
\text{minimum }=9\text{ reversal pairs}=18\text{ orbits}=108\text{ transitions},
\qquad
\#\{\text{optima}\}=2 .
\]

\subsection*{What changes and what does not}

\begin{center}
\begin{tabular}{lrrr}
\toprule
constraint & unit of choice & minimum transitions & optima \\
\midrule
none & single transition & 102 & --- \\
\(S_3\)-invariant & orbit (6 transitions) & 102 & 8 \\
\(S_3\) and reversal & reversal pair (12 transitions) & 108 & 2 \\
\bottomrule
\end{tabular}
\end{center}

Imposing reversal symmetry \emph{raises} the deletion cost from 102 to 108
while \emph{reducing} the number of independent choices from eight optima to
two. The two effects run in opposite directions and should not be conflated.

\section{Reproducibility and provenance}
\label{app:E}

This appendix records, for each machine-supported quantity, what artifact
supports it and at what provenance level. Provenance levels are used in a
deliberately restrictive sense:

\begin{description}
\item[\textsc{Package-local, re-executed}] the script ships in this package,
      reads no external file, and was re-run during preparation of this
      package; its output is included verbatim.
\item[\textsc{Reconstructed}] the script ships in this package, but no
      upstream original could be located, so byte-for-byte historical identity
      with the code that first produced the numbers cannot be established.
\item[\textsc{Copied, hash differs}] an upstream file of the same name was
      located, but its SHA-256 differs from the packaged copy; the packaged
      copy is therefore a variant, not a faithful archive.
\item[\textsc{Independent re-derivation}] the quantity was recomputed during
      refereeing by a separately written implementation.
\end{description}

\begin{center}
\begin{tabular}{p{0.36\textwidth}p{0.30\textwidth}p{0.26\textwidth}}
\toprule
Quantity & Artifact & Provenance level \\
\midrule
Boundary counts, \(|E(H)|\), \(\nu(H)=\tau(H)=102\), restricted pool
decomposition, \(S_3\) and reversal optima (Appendices~\ref{app:C},
\ref{app:D})
& \texttt{verify\_scissor\_and\_symmetry.js}; output in
\texttt{VERIFIER\_OUTPUT.txt}; certificates in the two JSON files
& \textsc{Package-local, re-executed} \\
\addlinespace
Period and lattice certificates, \(h=2,\ldots,6\) (Appendix~\ref{app:A})
& walks printed in full in Appendix~\ref{app:A}
& \textsc{Independent re-derivation}; every walk re-checked edge by edge \\
\addlinespace
Raw, valid, essential state counts (Table~\ref{tab:counts})
& \texttt{evidence/sft-container.js}
& \textsc{Copied, hash differs} \\
\addlinespace
Perron roots \(\lambda_h\) and variance rates \(a_h\) (Table~\ref{tab:ah})
& \texttt{evidence/spectral\_certifier.js}
& \textsc{Reconstructed}; values independently re-derived by two methods \\
\addlinespace
\(K=5\) ablation constants (Table~\ref{tab:ablation})
& \texttt{evidence/ablation\_certifier.js}
& \textsc{Reconstructed} \\
\addlinespace
Finite-\(n\) controls for \(h=2,3,4,5\) (Table~\ref{tab:h5})
& \texttt{evidence/independent\_dp.js}
& \textsc{Reconstructed} \\
\addlinespace
Finite-\(n\) data for \(h=6\) (Table~\ref{tab:h6}) and the correction audit
& \texttt{evidence/exact\_h6\_collision\_dp.js}
& \textsc{Copied, hash differs}; all five denominators
\textsc{Independent re-derivation} \\
\bottomrule
\end{tabular}
\end{center}

\subsection*{Hash record}

\begin{center}
\begin{tabular}{lll}
\toprule
file & packaged SHA-256 (first 16) & upstream SHA-256 (first 16) \\
\midrule
\texttt{sft-container.js} & \texttt{def436906535ca12} & \texttt{6dcf2a08252eba0b} \\
\texttt{exact\_h6\_collision\_dp.js} & \texttt{cada95990b11f3e9} & \texttt{bcaf8eabff09877e} \\
\texttt{spectral\_certifier.js} & \texttt{dfccc7c75ebf506c} & not located \\
\texttt{ablation\_certifier.js} & \texttt{3f7e8aa9530187c7} & not located \\
\texttt{independent\_dp.js} & \texttt{6802da24fa551d31} & not located \\
\bottomrule
\end{tabular}
\end{center}

For the two files with a located upstream, the hashes differ; we therefore do
not assert that the packaged copy is the code that originally produced the
reported numbers. For the three files with no located upstream, no historical
provenance claim is made at all. This is why the load-bearing quantities of
Appendices~\ref{app:A}, \ref{app:C} and~\ref{app:D}, and all five \(h=6\)
denominators, were re-derived independently rather than inherited.

\subsection*{Reporting incident}

The \(h=6\) reporting-only defect described in Appendix~\ref{app:B} is
retained here rather than silently repaired: the earlier prose report quoted
absolute counts that did not correspond to the preserved raw output, while the
\(nR_n\) ratios were unaffected. The correction is recorded, not overwritten.

\section{Local limit theorem details}
\label{app:F}

Throughout, \(M_h(\theta)\) denotes the Fourier-tilted transfer matrix of the
synchronized product chain on \(E_h\times E_h\): for a product edge
\(E=\bigl((x,y)\to(x',y')\bigr)\) carrying the rank-two increment
\(g(E)=\eta(y\to y')-\eta(x'\to x)\) of Section~\ref{sec:fold}, we set
\(M_h(\theta)_{(x,y),(x',y')}=e^{i\langle\theta,g(E)\rangle}\) when \(E\) is an
edge and \(0\) otherwise. Thus \(M_h(0)\) is the adjacency matrix of the
product graph, which is primitive by Proposition~\ref{prop:graph}(5); write
\(\varrho=\rho(M_h(0))\).

\subsection{Wielandt equality and the lattice obstruction}

\begin{lemma}
\label{lem:F1}
For every \(\theta\in[-\pi,\pi)^2\) with \(\theta\neq0\) we have
\(\rho(M_h(\theta))<\varrho\).
\end{lemma}

\begin{proof}
Entrywise \(|M_h(\theta)|=M_h(0)\), so \(\rho(M_h(\theta))\le\varrho\).
Suppose equality holds. Since \(M_h(0)\) is primitive, Wielandt's theorem
yields \(\varphi\in\mathbb R\) and a diagonal unitary \(D\) with
\(M_h(\theta)=e^{i\varphi}DM_h(0)D^{-1}\); equivalently there is
\(\psi:E_h\times E_h\to\mathbb R/2\pi\mathbb Z\) with
\[
\langle\theta,g(E)\rangle\equiv\varphi+\psi(u')-\psi(u)\pmod{2\pi}
\]
for every product edge \(E=(u\to u')\). Summing along a closed product walk
\(W\) of length \(L\), the coboundary telescopes, leaving
\begin{equation}\label{eq:F1cocycle}
\Bigl\langle\theta,\ \textstyle\sum_{E\in W}g(E)\Bigr\rangle\equiv L\varphi
\pmod{2\pi}.
\end{equation}

Let \(c_1,c_2,c_3\) be the three equal-length closed walks of length \(L\) at
the common root \(z_0\) recorded for this \(h\) in Appendix~\ref{app:A}. For
each pair \((i,j)\) let \(W_{ij}\) be the closed product walk at
\((z_0,z_0)\) that traverses \(c_i\) backwards in the first coordinate and
\(c_j\) forwards in the second. Then \(W_{ij}\) has length \(L\) and, by the
definition of \(g\),
\[
\sum_{E\in W_{ij}}g(E)=\Psi(c_j)-\Psi(c_i),
\]
the projected Parikh difference. Taking \(i=j\) makes the left-hand side of
\eqref{eq:F1cocycle} zero, whence \(L\varphi\equiv0\pmod{2\pi}\); the constant
\(\varphi\) is thereby eliminated with no separate argument. Substituting this
back into \eqref{eq:F1cocycle} for general \(i,j\) gives
\[
\bigl\langle\theta,\ \Psi(c_j)-\Psi(c_i)\bigr\rangle\equiv0\pmod{2\pi}.
\]
In particular \(\langle\theta,d_1\rangle\equiv\langle\theta,d_2\rangle\equiv0\)
for \(d_1=\Psi(c_2)-\Psi(c_1)\) and \(d_2=\Psi(c_3)-\Psi(c_1)\).

This is the only point at which the covolume of the certificate is used.
Appendix~\ref{app:A} records \(\det[d_1\ d_2]=\pm1\), so \(\{d_1,d_2\}\) is a
\(\mathbb Z\)-basis of \(\mathbb Z^2\); had the determinant been some \(q\)
with \(|q|>1\), the conclusion below would only give
\(\theta\in(2\pi/q)\mathbb Z^2\) and nontrivial characters would survive.
Unimodularity gives \(\langle\theta,v\rangle\equiv0\pmod{2\pi}\) for every
\(v\in\mathbb Z^2\), hence for \(v=e_1,e_2\), so \(\theta\in2\pi\mathbb Z^2\),
i.e.\ \(\theta=0\) in \([-\pi,\pi)^2\) --- a contradiction.
\end{proof}

\subsection{Uniform spectral decay away from \texorpdfstring{\(\theta=0\)}{theta=0}}

\begin{lemma}
\label{lem:F2}
For every \(\delta>0\) there exist \(\lambda_\delta<\varrho\) and
\(Q<\infty\) with
\[
\|M_h(\theta)^N\|\le Q\,\lambda_\delta^{\,N}
\qquad\text{for all }N\ge0\text{ and all }\theta\in K_\delta
:=\{\delta\le\|\theta\|_\infty\le\pi\}.
\]
\end{lemma}

\begin{proof}
The entries of \(M_h(\theta)\) are trigonometric polynomials, so
\(\theta\mapsto M_h(\theta)\) is continuous; in fixed finite dimension the
roots of the characteristic polynomial depend continuously on the entries, so
\(\theta\mapsto\rho(M_h(\theta))\) is continuous. By Lemma~\ref{lem:F1} it is
strictly less than \(\varrho\) at every point of the compact set
\(K_\delta\), hence
\[
\lambda'_\delta:=\max_{\theta\in K_\delta}\rho(M_h(\theta))<\varrho .
\]
A pointwise Gelfand estimate would not suffice here, since its implied
constant depends on \(\theta\). Fix \(\lambda_\delta\) with
\(\lambda'_\delta<\lambda_\delta<\varrho\) and let
\(C_\delta=\{z:|z|=\lambda_\delta\}\). For
\((z,\theta)\in C_\delta\times K_\delta\) we have
\(z\notin\operatorname{spec}(M_h(\theta))\), so the resolvent
\(R(z,\theta)=(zI-M_h(\theta))^{-1}\) exists; matrix inversion is continuous
where defined, so \((z,\theta)\mapsto\|R(z,\theta)\|\) is continuous on the
compact set \(C_\delta\times K_\delta\) and
\[
Q_0:=\sup_{C_\delta\times K_\delta}\|R(z,\theta)\|<\infty .
\]
Since the whole spectrum of \(M_h(\theta)\) lies strictly inside
\(C_\delta\), the holomorphic functional calculus gives
\(M_h(\theta)^N=\frac{1}{2\pi i}\oint_{C_\delta}z^NR(z,\theta)\,dz\), whence
\[
\|M_h(\theta)^N\|\le\frac{1}{2\pi}\cdot2\pi\lambda_\delta\cdot
\lambda_\delta^{\,N}\,Q_0=Q\,\lambda_\delta^{\,N},
\qquad Q:=Q_0\lambda_\delta,
\]
uniformly on \(K_\delta\).
\end{proof}

\subsection{Uniformity in endpoints, bounded displacement and lengths \texorpdfstring{\(N+r\)}{N+r}}

\begin{proposition}
\label{prop:F3}
Fix \(r_0\in\mathbb N\) and a finite \(V\subset\mathbb Z^2\). Let
\(\mathcal N_N(\alpha,\omega;v)\) be the number of length-\(N\) product paths
from \(\alpha\) to \(\omega\) with \(\sum_E g(E)=v\), and let
\(\mathcal W_N(\alpha,\omega)=(M_h(0)^N)_{\alpha,\omega}\). Then
\[
\mathcal N_{N+r}(\alpha,\omega;v)
=\mathcal W_{N+r}(\alpha,\omega)\,
\frac{1}{2\pi N\sqrt{\det\Sigma_{D,h}}}\,\bigl(1+o(1)\bigr),
\]
where the \(o(1)\) is uniform simultaneously over all state pairs
\((\alpha,\omega)\), all \(v\in V\), and all integers \(|r|\le r_0\).
\end{proposition}

\begin{proof}
Fourier inversion on \(\mathbb Z^2\) is exact:
\[
\mathcal N_{N+r}(\alpha,\omega;v)=\frac1{(2\pi)^2}
\int_{[-\pi,\pi]^2}e^{-i\langle\theta,v\rangle}
\bigl(M_h(\theta)^{N+r}\bigr)_{\alpha,\omega}\,d\theta .
\]
Split at \(\|\theta\|_\infty=\delta\). On \(K_\delta\), Lemma~\ref{lem:F2}
bounds the integrand by \(Q\lambda_\delta^{N+r}\) with
\(\lambda_\delta<\varrho\); this bound involves neither \(\alpha,\omega\) nor
\(v\), so that contribution is \(O(\lambda_\delta^{N})=o(\varrho^{N}/N)\)
uniformly in all three parameters.

Near the origin, \(M_h(0)\) is primitive, so \(\varrho\) is a simple isolated
eigenvalue. By analytic perturbation theory there are \(\delta_0>0\), an
analytic simple eigenvalue \(\varrho(\theta)\) and an analytic rank-one
spectral projection \(\Pi(\theta)\) on \(\|\theta\|_\infty<\delta_0\), with the
remaining spectrum inside \(\{|z|\le\varrho''\}\) for some
\(\varrho''<\varrho\); hence
\(M_h(\theta)^{N+r}=\varrho(\theta)^{N+r}\Pi(\theta)+O((\varrho'')^{N+r})\)
uniformly. Write \(\varrho(\theta)=\varrho\,e^{p(\theta)}\) with \(p\)
analytic and \(p(0)=0\). By Lemma~\ref{lem:cov2} and the \(S_3\) letter
frequencies \((1/3,1/3,1/3)\) the drift vanishes, so \(\nabla p(0)=0\) and
\(p(\theta)=-\tfrac12\theta^{\!\top}\Sigma_{D,h}\theta+O(\|\theta\|^3)\).

Substituting \(\theta=\phi/\sqrt N\) and using \(\nabla p(0)=0\),
\[
\varrho(\theta)^{N}=\varrho^{N}
\exp\bigl(-\tfrac12\phi^{\!\top}\Sigma_{D,h}\phi\bigr)(1+o(1)),
\]
and \(\Pi(\theta)\to\Pi(0)\). Gaussian integration produces the factor
\((2\pi N)^{-1}(\det\Sigma_{D,h})^{-1/2}\).

The three uniformities are now read off. \emph{Endpoints:} there are finitely
many pairs \((\alpha,\omega)\) and the dependence enters only through the entry
\(\Pi(\theta)_{\alpha,\omega}\), so the maximum of finitely many \(o(1)\)
terms is \(o(1)\). \emph{Displacement:} \(v\) enters only through
\(e^{-i\langle\theta,v\rangle}=e^{-i\langle\phi,v\rangle/\sqrt N}\), and
\(\bigl|e^{-i\langle\phi,v\rangle/\sqrt N}-1\bigr|\le\|\phi\|\,\|v\|/\sqrt N\),
which is uniform over the fixed finite set \(V\). \emph{Length:}
\(\varrho(\theta)^{N+r}=\varrho(\theta)^{N}\varrho(\theta)^{r}\) and
\(\varrho(\theta)^{r}\to\varrho^{r}\) uniformly for \(|r|\le r_0\), while
\((2\pi(N+r))^{-1}=(2\pi N)^{-1}(1+O(1/N))\).
\end{proof}

\begin{remark}
This is exactly the form consumed by Lemma~\ref{lem:boundary}. After peeling,
at most \(D_h\) edges are removed at each end, so the surviving core portion
has length \(N-r\) with \(0\le r\le 2D_h\); its endpoints are the finitely many
core states at which the path enters and leaves \(E_h\); and by
Lemma~\ref{lem:fold} the required displacement is \(\beta\) adjusted by the
Parikh contribution of the peeled edges, so it lies in the fixed finite set
\(\{v:\|v\|_\infty\le m+2D_h\}\). The resulting relative error is therefore a
single \(o(1)\) independent of the term, and since every term of the identity
of Lemma~\ref{lem:pathid} is non-negative, summation preserves it.
\end{remark}


\end{document}
